\documentclass[11pt,a4paper]{article}
\usepackage[T1]{fontenc}
\usepackage[utf8]{inputenc}
\usepackage{mathptmx}
\AtBeginDocument{\let\hbar\relax
  \DeclareMathSymbol{\hbar}{\mathord}{AMSb}{"7E}}
\usepackage{amsmath,amssymb,amsthm}
\usepackage{geometry}
\usepackage{microtype}
\usepackage{hyperref}
\usepackage{booktabs}
\usepackage{enumitem}
\usepackage[numbers]{natbib}
\usepackage{titlesec}

\geometry{top=1.1in,bottom=1.1in,left=1.15in,right=1.15in}

\titleformat{\section}{\normalfont\bfseries\large}{\thesection}{1em}{}
\titleformat{\subsection}{\normalfont\bfseries\normalsize}{\thesubsection}{1em}{}
\titleformat{\subsubsection}{\normalfont\bfseries\itshape\normalsize}
  {\thesubsubsection}{1em}{}
\titlespacing*{\section}{0pt}{14pt}{4pt}
\titlespacing*{\subsection}{0pt}{10pt}{3pt}

\theoremstyle{plain}
\newtheorem{theorem}{Theorem}[section]
\newtheorem{proposition}[theorem]{Proposition}
\newtheorem{corollary}[theorem]{Corollary}
\theoremstyle{definition}
\newtheorem{definition}[theorem]{Definition}
\newtheorem{prediction}[theorem]{Prediction}
\theoremstyle{remark}
\newtheorem{remark}[theorem]{Remark}

\hypersetup{colorlinks=true,linkcolor=black,citecolor=black,urlcolor=black,
  pdftitle={Relational Actualism: A Biography of the Universe},
  pdfauthor={Joshua F. Sandeman}}

\newcommand{\Pact}{\mathcal{P}_{\mathrm{act}}}

% ─────────────────────────────────────────────────────────────────────────────
\begin{document}

% ── Title block ───────────────────────────────────────────────────────────────
\begin{center}
  {\LARGE\bfseries Relational Actualism:\par}
  \vspace{4pt}
  {\large\bfseries A Biography of the Universe\par}
  \vspace{6pt}
  {\normalsize Paper~0 of the Relational Actualism Series\par}
  \vspace{14pt}
  {\large Joshua F.\ Sandeman$^{1*}$\par}
  \vspace{4pt}
  {\normalsize $^{1}$Independent Researcher, Salem, OR, USA.\par}
  {\normalsize $^{*}$Corresponding author: sansuikyo@gmail.com
   $\cdot$ relationalactualism.org\par}
  \vspace{14pt}
\end{center}

\begin{quote}\itshape
The universe is not a thing that exists.  It is a thing that happens---and
it never stops happening.
\end{quote}

% ── Abstract ─────────────────────────────────────────────────────────────────
\begin{abstract}
This paper tells a single story: the life of a universe, from the dynamics
of the parent supermassive black hole through the emergence of particles,
spacetime, galaxies, and life, to a future of perpetual branching rather
than heat death.  The story is told in the language of Relational
Actualism~(RA), a framework built on a single ontological commitment: that
irreversible actualization events are the primitive physical reality.
From this commitment, implemented as an irreversibly growing directed
acyclic graph with a unique self-consistent growth rule, the framework
derives the coupling constants, the Standard Model spectrum and gauge
groups, the structural zero cosmological constant, the dark matter
mechanism, the Hubble tension resolution, five CMB anomalies as signatures
of parent SMBH dynamics, the dissolution of the measurement problem, and
substrate-independent conditions for biological complexity---all with no
free parameters.  The theory's logical structure is a constraint web, not
a deduction chain, and its mathematical core has been formally verified
(163~Lean~4 theorems, compiled clean).  This paper maps the complete
framework, catalogues its predictions, and presents the formal
verification record.  The companion papers~(P1--P4) provide the
derivations; this paper provides the meaning.
\end{abstract}

%================================================================
\section*{Preface}\label{sec:preface}
%================================================================

This document tells a single story: the life of a universe.  Not a
universe described by equations, but a universe that \emph{does}
something---that grows, selects, branches, and complexifies by its own
nature.  The story begins before the Big Bang, follows the universe from
birth through the emergence of galaxies, stars, chemistry, and life, and
looks forward to a future that is neither heat death nor eternal stasis,
but perpetual creative transformation.

The story is told in the language of Relational Actualism~(RA), a
framework built on a single ontological commitment: that irreversible
actualization events---physical interactions that produce permanent
increases in quantum relative entropy---are the primitive physical
reality.  They are not events in spacetime; spacetime is the macroscopic
description of their accumulated pattern.  From this one commitment, a
coherent account cascades across the full range of physics, from the
values of the coupling constants to the substrate-independent conditions
for biological complexity.

The technical machinery behind each piece of the story has been developed
across a suite of companion papers, each with its own mathematical
apparatus and formal verification.  This document does not reproduce those
derivations.  Instead, it tells the story those derivations support, in
the order the universe itself tells it.  The reader who wants the proof
that $\alpha_{\mathrm{EM}}^{-1} = 137.036$ from integer arithmetic, or
the Lean~4 verification of causal invariance, or the DFT calculations
supporting the Causal Firewall, will find them in the referenced papers.
Here, they will find what those results \emph{mean}---what they tell us
about the universe we inhabit, and why the fit matters.

One result deserves special mention.  The cosmic microwave
background---the oldest light in the universe---carries anomalies that
standard cosmology cannot explain: an aligned quadrupole and octupole,
a suppressed quadrupole power, and a hemisphere asymmetry.  Relational
Actualism identifies these as signatures of the spinning supermassive
black hole from which our universe nucleated.  Three parameters (the
parent's spin, the duration of inflation, and the accretion asymmetry)
account for five independent observations, with no free parameters
remaining.  We are reading the imprint of our origin in the sky.

%================================================================
\section{Act~I: Before the Beginning}\label{sec:act1}
%================================================================

Every universe, in the RA picture, is born from another.  Not as metaphor,
but as mechanism: when a region of causal graph reaches its maximum
actualization rate---the bandwidth limit of the local causal structure---it
undergoes a \emph{causal severance}, a topological disconnection that
creates a new, causally independent graph.  The daughter universe inherits
its initial conditions from the boundary of the severance and then grows
according to the same rules as its parent.  The rules do not change from
one generation to the next, because they are not imposed on the graph from
outside.  They \emph{are} the graph.

\subsection{The maximum frame rate}\label{sec:framerate}

Think of the causal graph as a network being built in real time.  Each
actualization event writes a new vertex and new directed edges into the
graph.  The rate at which vertices can be added is not unlimited.  The
Local Ledger Condition~(LLC), which enforces conservation at every vertex,
constrains how many events can be processed per unit of graph time.  There
is a maximum: the Planck rate, approximately $10^{43}$~events per Planck
time per Planck volume.  This is the universe's frame rate.

When a collapsing region approaches this maximum, something decisive
happens.  The actualization density~$\mu$ climbs toward the causal
termination threshold $\mu_2 \approx 5.7$, where the acceptance
probability for new actualization events drops to zero.  The graph cannot
grow any further locally.  The region severs from its parent: no outgoing
edges remain, and the graph cut theorem (machine-verified in Lean~4)
guarantees that the LLC holds independently on each side.  The daughter
graph is born.

\subsection{Why a supermassive black hole}\label{sec:smbh}

Not every black hole can birth a viable universe.  The daughter inherits
its baryon-to-photon ratio~$\eta_b$ from the Bekenstein--Hawking entropy
of the parent at the moment of severance.  Our universe has
$\eta_b = 6.1 \times 10^{-10}$, which requires a parent entropy of
approximately $10^{90}\,k_B$.  Inverting the Bekenstein--Hawking formula
gives the mass of the parent: approximately 2.5 to 6.6 million solar
masses---a supermassive black hole, the kind found at the centres of
galaxies.

\subsection{The parent's geometry}\label{sec:kerr}

A real astrophysical SMBH spins, and its spin defines a preferred axis and
a preferred plane.  The Kerr metric describes a horizon that is oblate:
wider at the equator, compressed at the poles.  The degree of oblateness
is controlled by the dimensionless spin parameter~$a_*$.  Observed SMBHs
in the relevant mass range typically have $a_* \in [0.7, 0.95]$.

The parent also has an accretion environment: a disk of infalling matter
in its equatorial plane, with jets along its spin axis.  This environment
is not perfectly symmetric about the equatorial plane.  These
imperfections, as we will see, are preserved in the daughter's initial
conditions and are visible in our sky today.

%================================================================
\section{Act~II: Nucleation}\label{sec:act2}
%================================================================

The moment of severance is not an explosion.  It is a
disconnection---the last outgoing causal edge is written, and then there
are no more.  The daughter graph begins with the boundary conditions
inherited from the severance surface.

\subsection{The angular imprint}\label{sec:imprint}

The Kerr horizon's geometry imprints a specific angular pattern on the
boundary flux.  Decomposing this pattern into spherical harmonics reveals
a clean structure: the spin produces a pure quadrupole ($\ell = 2$)
deformation, with amplitude~$\varepsilon_2$ that grows with the spin
parameter.  The octupole ($\ell = 3$) is identically zero from the Kerr
geometry alone, because the Kerr metric is symmetric about its equatorial
plane.  Any octupole must come from the accretion environment---the
asymmetric infall that breaks the equatorial symmetry.

Both the quadrupole and the octupole are aligned with the same axis: the
parent's spin.  This is not a coincidence.  The quadrupole is set by the
spin geometry; the octupole is set by the accretion disk, which lies in
the spin's equatorial plane.  They \emph{must} share an axis.

\subsection{Inflationary dilution}\label{sec:inflation}

Immediately after nucleation, the daughter universe undergoes rapid
expansion.  This dilutes the nucleation perturbation exponentially.  A
mode with angular scale~$\ell$ is diluted by a factor $(2/\ell)^2$
relative to the largest mode, because shorter-wavelength modes exit the
causal horizon later and spend more time being stretched.

For the signal to be visible in the CMB, the universe must have undergone
\emph{nearly the minimum} amount of inflation needed to solve the horizon
problem---approximately 64 to 65~e-folds.  The dilution law predicts a
specific hierarchy: strong at $\ell = 2$, moderate at $\ell = 3$,
negligible by $\ell = 4$.

%================================================================
\section{Act~III: The First Moments}\label{sec:act3}
%================================================================

\subsection{The rules crystallise}\label{sec:rules}

As the daughter universe expands and cools, the causal graph's local
structure determines which patterns of actualization are stable.  The BDG
action---the natural action functional for a four-dimensional causal set,
characterised by the integers $(1,-1,9,-16,8)$---admits exactly five
topologically distinct stability classes, corresponding precisely to the
particle types of the Standard Model: quarks, gluons, gauge bosons, the
Higgs, and leptons.

The coupling constants emerge from the same integers.
$\alpha_{\mathrm{EM}}^{-1} = 137.036$ (agreement: $0.00001\%$).
$\alpha_s(m_Z) = 1/\sqrt{72} \approx 0.11785$ (agreement: $0.13\%$).
The Koide lepton mass formula $K = 2/3$.  The gauge group
$\mathrm{SU}(3)\times\mathrm{SU}(2)\times\mathrm{U}(1)$, determined by
the alternating signs of the BDG coefficients via the
inclusion--exclusion principle.  None of these are inputs.  They are
outputs of four-dimensional causal geometry.  The full derivations are
given in Paper~I~\cite{P1}.

\subsection{The selectivity of the universe}\label{sec:selectivity}

The actualization threshold
$\Delta S^* = 0.601$~nats---the minimum increase in quantum relative
entropy required for an interaction to become a permanent event---is
computed exactly from the BDG dynamics at the Planck density $\mu = 1$.
This threshold governs quantum measurement, the transition from quantum
to classical behaviour, and the origin of biological complexity.  A
single number, derived from four-dimensional geometry, that controls
physics at every scale.

%================================================================
\section{Act~IV: Structure}\label{sec:act4}
%================================================================

\subsection{The cosmic microwave background}\label{sec:cmb}

Approximately 380,000~years after nucleation, neutral atoms formed and
the CMB photons began their journey.  Standard cosmology predicts
statistically isotropic perturbations.  At small angular scales
($\ell \geq 4$), this prediction holds beautifully.  At the largest
scales, several anomalies have persisted through two decades of
measurement:

\emph{The low quadrupole.}  $C_2 \approx 200\,\mu\mathrm{K}^2$, versus
the $\Lambda$CDM prediction of $\sim\!1100\,\mu\mathrm{K}^2$---an 80\%
suppression.

\emph{The axis of evil.}  The $\ell = 2$ and $\ell = 3$ preferred axes
are aligned to within $9$--$13^\circ$, with no known mechanism in
standard cosmology.

\emph{The hemisphere asymmetry.}  The northern hemisphere (relative to
the preferred axis) has systematically less power than the southern.

\emph{Normal higher multipoles.}  $\ell \geq 4$ modes match $\Lambda$CDM.

RA explains all five features with a single mechanism.  The parent's spin
produces a coherent quadrupole that partially cancels the stochastic
component, suppressing~$C_2$.  The accretion asymmetry produces a
coherent octupole aligned with the same axis.  The hemisphere asymmetry
is the odd-parity $P_3$ component.  The $(2/\ell)^2$ dilution ensures
the signal is negligible above $\ell = 3$.  Three parameters ($a_*$,
$\Delta N$, $\delta_{\mathrm{acc}}$), three constraints, two consistency
checks, zero remaining freedom.  The full treatment is in
Paper~III~\cite{P3}.

\subsection{The cosmic web}\label{sec:web}

The RA field equations couple the local expansion rate to baryon density:
lower density expands faster (further diluting); higher density expands
slower (further concentrating).  This positive feedback drives
filamentary cosmic web topology as a dynamical attractor.

Dark matter, in this picture, is not a particle species.  It is the
actualization-density topology of sparse graph regions, where effective
dimensionality reduces from four to two, producing flat rotation curves
from baryonic matter alone.  The WIMP prohibition follows structurally:
any particle without on-shell electromagnetic or strong interactions has
zero actualization depth and zero gravitational coupling.

\subsection{The Hubble tension}\label{sec:hubble}

The 4--5$\sigma$ disagreement in $H_0$ measurements is a natural
consequence of density-dependent expansion.  The predicted value for the
Eridanus supervoid is $H_0 \approx 75.9$~km/s/Mpc, testable with DESI:
RA predicts a linear correlation between $H_0$ and baryon density along
the line of sight.

%================================================================
\section{Act~V: Complexity}\label{sec:act5}
%================================================================

The same actualization filter, operating recursively at successively
larger scales, produces the full hierarchy of physical complexity.  The
Erd\H{o}s--R\'enyi percolation threshold at $\mu = 1$, which governs QCD
confinement and filamentary cosmic structure, also governs the transition
from abiotic chemistry to biological complexity.

\subsection{The Causal Firewall}\label{sec:firewall}

Life is possible wherever two conditions are met: redundant information
recording~(F1) and controlled actualization rate~(F2).  These are
substrate-independent---they make no reference to carbon, water, or
molecular chirality.  They imply an origin-of-life sandwich bound and
universal biosignature criteria.  The full development is in
Paper~IV~\cite{P4}.

\subsection{Complexification and dissipation}\label{sec:complex}

Every actualization event is simultaneously an increase in entropy
\emph{and} a commitment to a specific structure.  Dissipation and
complexification are not opposing forces.  They are two faces of the same
act.  The universe builds complexity \emph{by} producing entropy, because
the actualization filter selects for structured outcomes at every scale.

%================================================================
\section{Act~VI: The Quantum World}\label{sec:act6}
%================================================================

\subsection{The measurement problem dissolved}\label{sec:measurement}

A quantum state is any state with relative entropy below $\Delta S^*$.
When environmental interaction pushes the relative entropy past the
threshold, a new permanent vertex is written.  Nothing ``collapses.''
There is no Heisenberg cut, because there is no fundamental distinction
between quantum and classical---only a gradient of actualization density.
The full treatment, including the Born rule consistency proof and the
causal invariance theorem (machine-verified in Lean~4), is in
Paper~II~\cite{P2}.

\subsection{General relativity: inevitable}\label{sec:gr}

Einstein's field equations are not imposed on RA.  They are derived from
it.  The Local Ledger Condition~(L01, Lean-verified) gives local
conservation at every vertex.  Amplitude locality~(O01, Lean-verified)
and $d = 4$~(L11, Lean-verified) together with the Benincasa--Dowker
uniqueness theorem~\cite{BenincasaDowker2010} determine the BDG action
as the unique local causal-set action, whose continuum limit is the
Einstein--Hilbert action~\cite{BenincasaDowker2010}.  The Bianchi
identity is not a separate geometric fact---it is the LLC viewed at the
level of the metric.  The actualization projector removes vacuum
contributions, giving $\Lambda = 0$ exactly.  General relativity is not a
fundamental theory in RA.  It is the effective description of causal graph
dynamics at scales far above the Planck length.

\subsection{Quantum computing: a hard ceiling}\label{sec:qc}

Because each qubit is an additional kinematic site at which actualization
events can occur, RA predicts a structural ceiling on fault-tolerant
quantum array size: $N_{\max} = \eta \cdot p_{\mathrm{th}}$.  Beyond this
size, spurious actualizations exceed the error budget regardless of
hardware quality.  Testable at the scale of hundreds to thousands of
logical qubits.

%================================================================
\section{Act~VII: The Fate of the Universe}\label{sec:act7}
%================================================================

\subsection{Fragmentation, not heat death}\label{sec:fragmentation}

The void--filament feedback does not stop.  Voids approach the
percolation threshold $\mu = 1$, below which the graph fragments
quietly.  Dense boundaries approach $\mu_2 \approx 5.7$, where new
severances fire and daughter universes nucleate.  The universe does not
wind down.  It branches.

This also prohibits Boltzmann Brains: no causal patch persists long
enough for the infinite-time limit that Boltzmann Brain arguments require.

\subsection{The cycle}\label{sec:cycle}

At the dense boundaries, the cycle begins again.  The daughter inherits
its initial conditions from the parent's local structure.  Angular
momentum is transmitted across generations, encoded in the boundary
conditions of each successive severance.  The rules do not change from
generation to generation, because they were never imposed from outside.

%================================================================
\section*{Coda: One Nature of Nature}\label{sec:coda}
%================================================================

The story told in this document involves no new fields, no extra
dimensions, no additional free parameters, and no post-hoc adjustments.
Every piece follows from the same ontological commitment: that
irreversible actualization events, writing permanent edges into a growing
directed acyclic graph, are the primitive physical reality.

The deepest insight of Relational Actualism is that the traditional
distinction between ``laws'' and ``phenomena'' is an artifact of
description.  The Local Ledger Condition is not a law that the causal
graph obeys.  It \emph{is} the causal graph, stated as a property.
Conservation of energy is not a constraint imposed on actualization
events.  It is what actualization events \emph{are}, viewed at the level
of bookkeeping.  The irreversibility of time is not explained by the
framework.  It is \emph{encoded} in the framework's most primitive
commitment.

The web of mutual implications revealed by the dependency structure of
the RA programme is not a feature of the theory.  It is a feature of
\emph{Nature} that the theory makes visible.

\begin{quote}\itshape
The universe is not a thing that winds down.  It is a thing that
branches---endlessly, irreversibly, creatively.  It has no terminal state,
because the process that builds complexity is the same process that
fragments the graph and nucleates new worlds.  Complexification and
dissipation are two faces of the same act: ceaseless, causal
transformation.  This is what one nature of Nature looks like, from bosons
to brains and beyond.
\end{quote}

%================================================================
\section{The Constraint Web}\label{sec:constraintweb}
%================================================================

The logical structure of RA is not a deduction chain but a
\emph{constraint web}: a network of mutually reinforcing results, each
derivable from multiple independent paths.  Approximately nine
interlocking loops have been identified:

\begin{enumerate}[nosep,leftmargin=2em]
\item \textbf{BDG $\to$ couplings $\to$ spectrum.}
  The same integers determine coupling constants and particle types.
\item \textbf{LLC $\to$ GR.}
  Conservation at every vertex, via BDG uniqueness, becomes the Einstein
  field equations in the continuum approximation.
\item \textbf{$\Delta S^* \to$ measurement $\to$ Born rule.}
  The actualization threshold governs quantum measurement.
\item \textbf{$\mu = 1 \to$ five scales.}
  The percolation threshold appears in QCD, galactic dynamics, quantum
  computing, selectivity, and biology.
\item \textbf{$\Pact \to \Lambda = 0$, WIMP prohibition.}
  The actualization projector removes vacuum contributions and excludes
  non-interacting particles from gravity.
\item \textbf{Severance $\to$ nucleation $\to$ CMB.}
  Parent Kerr geometry imprints coherent multipoles on the daughter.
\item \textbf{Uniqueness.}
  The BDG coefficients are the unique self-consistent solution.
\item \textbf{Void--filament $\to$ fragmentation.}
  Density feedback drives cosmic web topology and prohibits heat death.
\item \textbf{F1~+~F2 $\to$ Causal Firewall.}
  Substrate-independent conditions for biological complexity.
\end{enumerate}

The web's self-consistency provides combinatorial evidence for
correctness: the probability that all connections are simultaneously wrong
by coincidence decreases with the number of independent paths.

%================================================================
\section{Formal Verification Record}\label{sec:lean}
%================================================================

The discrete mathematical core has been formally verified in
Lean~4.29/Mathlib.

\begin{table}[ht]
\centering\small
\begin{tabular}{@{}lrrl@{}}
\toprule
File & Theorems & Sorry & Key results \\
\midrule
\texttt{RA\_D1\_Proofs}       & 76 & 0 & $\alpha_{\mathrm{EM}}^{-1}=137$, Koide, SM spectrum \\
\texttt{RA\_GraphCore}         &  5 & 0 & LLC, Graph Cut, Markov Blanket \\
\texttt{RA\_Koide}             & 16 & 0 & $K = 2/3$ for all $\theta$, $m_0$ \\
\texttt{RA\_AmpLocality}       &  5 & 0 & Amplitude locality, causal invariance \\
\texttt{RA\_AQFT\_Proofs\_v10} &  9 & 1 & Frame independence, Rindler \\
\texttt{RA\_O14\_Uniqueness}   & 50 & 0 & BDG uniqueness (Yeats--Rota) \\
\texttt{RA\_Spin2\_Macro}      &  2 & 0 & Massless spin-2 DOF count \\
\midrule
\textbf{Total} & \textbf{163} & \textbf{1} & \\
\bottomrule
\end{tabular}
\caption{Lean~4 verification record.  The single \texttt{sorry} requires
an external library dependency; it is not a conceptual gap.}
\label{tab:lean}
\end{table}

%================================================================
\section{Predictions}\label{sec:predictions}
%================================================================

The following predictions are falsifiable.  A single confirmed failure
would refute the framework.

\begin{description}[style=nextline,leftmargin=1.5em,nosep]
\item[Near-term]
  BMV null result (gravity-mediated entanglement);
  $H_0$ gradient with baryon density (DESI);
  no WIMP signal below neutrino floor;
  spin-bath collapse timescale $t^* = 0.274/g$.
\item[Medium-term]
  Hard ceiling on fault-tolerant quantum arrays ($N_{\max}$);
  CMB axis persistence in CMB-S4 at ${>}\,3\sigma$;
  neutrino mass sum $\sum m_\nu \approx 59$~meV;
  cosmic web orientation correlates with CMB axis.
\item[Long-term]
  Causal Firewall biosignatures (substrate-independent);
  no heat death (fragmentation precedes equilibrium).
\end{description}

%================================================================
\section{Relation to Other Programmes}\label{sec:other}
%================================================================

\textbf{Causal set theory.}  RA is built on causal set
foundations~\cite{BLMS1987,Sorkin2005} and extends the programme by
deriving the BDG coefficients intrinsically and extracting physical
observables from them.

\textbf{Loop quantum gravity.}  LQG discretises the state space; RA
discretises the growth process.  RA makes specific predictions (coupling
constants, particle spectrum) that LQG does not.

\textbf{String theory.}  No extra dimensions, no supersymmetry, no
landscape.  Opposite BSM predictions.

\textbf{Assembly Theory.}  RA's assembly depth $\tilde{A}_{\mathrm{RA}}$
is grounded in fundamental physics; Assembly Theory's index in chemical
operations.  Paper~IV shows convergence for complex molecules.

%================================================================
\section{Open Problems}\label{sec:open}
%================================================================

Zero foundational problems remain open.  Five technical problems await
resolution: WEP formal bounds~(O04), Unruh detector model~(O05), BMV
timescale~(O07, \emph{highest priority}), Causal Firewall
$\tau_d$~(O08), and fragmentation timescale~(O13).

% ─────────────────────────────────────────────────────────────────────────────
\section*{The Companion Papers}
% ─────────────────────────────────────────────────────────────────────────────

\textbf{P1}---\emph{Fine Structure Constant and Strong Coupling from
Four-Dimensional Causal Geometry.}  Derives $\alpha_{\mathrm{EM}}$ and
$\alpha_s$ from BDG integers.  The rules crystallise (Act~III).

\textbf{P2}---\emph{Wave Function Collapse as Irreversible Entropy
Production.}  Dissolves the measurement problem.  The quantum world
(Act~VI).

\textbf{P3}---\emph{The Origin, Structure, and Fate of Universes.}
Nucleation cosmology, CMB anomalies, dark matter, Hubble tension, cosmic
web, fragmentation.  Acts~I, II, IV, and~VII.

\textbf{P4}---\emph{The Causal Firewall.}  Substrate-independent
conditions for biological complexity.  Act~V.

\medskip\noindent
Lean~4 proof files and computation scripts are available at
\url{github.com/jsandeman/Relational-Actualism}.  All papers are
open-access on Zenodo.  Correspondence: sansuikyo@gmail.com.

% ─────────────────────────────────────────────────────────────────────────────
\begin{thebibliography}{99}

\bibitem{P1}
J.~F.~Sandeman,
``Paper~I: Fine Structure Constant and Strong Coupling from
Four-Dimensional Causal Geometry,''
Preprint (2026).

\bibitem{P2}
J.~F.~Sandeman,
``Paper~II: Wave Function Collapse as Irreversible Entropy Production,''
Preprint (2026).

\bibitem{P3}
J.~F.~Sandeman,
``Paper~III: The Origin, Structure, and Fate of Universes,''
Preprint (2026).

\bibitem{P4}
J.~F.~Sandeman,
``Paper~IV: The Causal Firewall,''
Preprint (2026).

\bibitem{BenincasaDowker2010}
D.~M.~T.~Benincasa and F.~Dowker,
``The scalar curvature of a causal set,''
\textit{Phys.\ Rev.\ Lett.}\ \textbf{104}, 181301 (2010).

\bibitem{DowkerGlaser2013}
F.~Dowker and L.~Glaser,
``Causal set d'Alembertians for various dimensions,''
\textit{Class.\ Quantum Grav.}\ \textbf{30}, 195016 (2013).

\bibitem{BLMS1987}
L.~Bombelli, J.~Lee, D.~Meyer, and R.~D.~Sorkin,
``Space-time as a causal set,''
\textit{Phys.\ Rev.\ Lett.}\ \textbf{59}, 521 (1987).

\bibitem{Sorkin2005}
R.~D.~Sorkin,
``Causal sets: Discrete gravity,''
in \textit{Lectures on Quantum Gravity}
(Springer, 2005), pp.~305--327.

\bibitem{Sharma2023}
M.~D.~Sharma \textit{et al.},
``Assembly theory explains and quantifies selection and evolution,''
\textit{Nature} \textbf{622}, 278--283 (2023).

\end{thebibliography}

\end{document}
